% Part: computability
% Chapter: computability-theory
% Section: complete-ce-sets

\documentclass[../../../include/open-logic-section]{subfiles}

\begin{document}

\olfileid{cmp}{thy}{cce}
\olsection{సంపూర్ణ గణనీయంగా లెక్కించదగిన సమితులు}

\begin{defn}
కింది షరతులు నెరవేరితే సమితి $A$ను (అనేకం-ఒకటి తగ్గించదగినత కింద)
\emph{సంపూర్ణ \tetoken{గణనీయంగా లెక్కించదగిన}{!!{computably enumerable}} సమితి} అంటాం:
\begin{enumerate}
\item $A$ గణనీయంగా లెక్కించదగినది; మరియు
\item గణనీయంగా లెక్కించదగిన ప్రతి ఇతర సమితి $B$కు $B\leq_m A$.
\end{enumerate}
\end{defn}

మరో మాటలో, సంపూర్ణ గణనీయంగా లెక్కించదగిన సమితులు, సాధ్యమైన
గణనీయంగా లెక్కించదగిన సమితులన్నిటిలోనూ ``అత్యంత కష్టమైనవి.''
\emph{ఏ} గణనీయంగా లెక్కించదగిన సమితి గురించిన ప్రశ్నలకైనా సమాధానం
ఇవ్వడానికి అవి తోడ్పడతాయి.

\begin{thm}
$K$, $K_0$, $K_1$ అన్నీ సంపూర్ణ గణనీయంగా లెక్కించదగిన సమితులు.
\end{thm}

\begin{proof}
$K_0$ సంపూర్ణమని చూడటానికి, $B$ ఏ గణనీయంగా లెక్కించదగిన సమితైనా
అనుకుందాం. అప్పుడు ఏదో సూచిక~$e$కు
\[
B=W_e=\Setabs{x}{\cfind{e}(x)\fdefined}.
\]
$f(x)=\tuple{e,x}$గా ప్రమేయం~$f$ను నిర్వచించండి. అప్పుడు ప్రతి సహజ
సంఖ్య~$x$కు $x\in B$ అయితేనే $f(x)\in K_0$. మరో మాటలో, $f$ అనేది
$B$ను~$K_0$కు తగ్గిస్తుంది.

$K_1$ సంపూర్ణమని చూడటానికి, \olref[k1]{prop:k1} నిరూపణలో $K_0$ను
దానికి తగ్గించామని గమనించండి. కనుక \olref[ppr]{prop:trans-red}
ప్రకారం గణనీయంగా లెక్కించదగిన ప్రతి సమితినీ~$K_1$కు తగ్గించవచ్చు.

దాదాపు అదే విధంగా $K_0$ను~$K$కు తగ్గించవచ్చు.
\sourcecorrection{OLTECOMTHY-015}{ఇక్కడ $K$ సంపూర్ణమని నిరూపించడానికి $K_0\leq_m K$ కావాలి; కానీ ఆంగ్ల మూలం ఇప్పటికే తెలిసిన వ్యతిరేక దిశ $K\leq_m K_0$ను రాసింది. సిద్ధాంతపు వాదనకు అవసరమైన దిశగా సరిచేశాం; కింది సమస్యలోనూ అదే దిశను సరిచేశాం.}
\end{proof}

\begin{prob}
$K_0$ను $K$కు తగ్గించే ప్రమేయాన్ని ఇవ్వండి.
\end{prob}

\begin{digress}
ఇప్పటివరకు మనం పరిశీలించిన గణనీయంగా లెక్కించదగిన సమితుల ఉదాహరణలన్నీ
గణనీయాలు లేదా సంపూర్ణాలు అని తేలింది. ఇది వింతగా అనిపించాలి!{}
గణనీయమూ కాని, సంపూర్ణమూ కాని గణనీయంగా లెక్కించదగిన సమితులు ఏమైనా
ఉన్నాయా? సమాధానం అవును; కానీ 1950ల మధ్యకాలంలోనే ఫ్రీడ్బర్గ్,
ముచ్నిక్ ఒకరికొకరు స్వతంత్రంగా దీన్ని స్థాపించారు.
\end{digress}

\end{document}
